%-----------------------------------------------------------------------------
\chapter{\babDua}

%------------------------------------------------------------

\section{Rice}
\label{sec:rice}

Rice (\textit{Oryza sativa} L.) is a staple food for more than half of the world's population, with Asia accounting for the dominant share of global production and consumption \cite{mohidem2022rice}. In Indonesia, the 2023 Agricultural Census reports that 17{,}251{,}432 of 27{,}802{,}434 agricultural land users, about 62 percent, are smallholder farmers cultivating less than 0.5 hectares, and food crops including rice form the dominant subsector with 15{,}550{,}786 farming households \cite{bps2023sensus}. Any technological intervention for rice productivity must therefore remain affordable, easy to operate, and deployable on sub-hectare plots that fall well below the scale of industrial precision agriculture.

The rice life cycle is conventionally divided into three growth phases, vegetative, reproductive, and ripening, with a complete cycle of 70 to 170 days depending on cultivar \cite{mohapatra2025staging}. The vegetative phase covers germination, tillering, and leaf development, during which biomass accumulates and the canopy closes. The reproductive phase, also called the generative phase, starts with panicle initiation and continues through booting, heading, and flowering. The ripening phase covers grain filling and maturation, during which starch is translocated to the grains and the canopy senesces toward harvest. This study uses the simplified labels vegetative, generative, and mature, where mature corresponds to the ripening phase and denotes the harvest-ready state.

The transition into the mature phase is the most operationally critical because it determines the harvest window. A meta-analysis of harvest-time effects on rice yield and quality shows that early harvesting reduces yield and produces chalky immature grains, while delayed harvesting in tropical regions lowers the head rice rate and increases shattering losses \cite{zhou2025metaharvest}. Phase boundaries are not rigid in calendar time. Controlled experiments on three rice varieties showed that elevated temperature significantly alters phenological intervals and reduces yield even within the same cultivar \cite{sanwong2023ricetemp}, and China's seasonal rice calendar exhibited measurable shifts in transplanting, heading, and maturity dates between 2003 and 2022 \cite{li2024chinaricecalendar}. A fixed Days After Sowing schedule cannot reliably indicate the harvest moment for every plot, and visual judgment of phase transitions, while accessible to smallholder farmers, is subjective and prone to inter-observer variability.

These characteristics motivate the technical scope of this study. Classifying phenology into vegetative, generative, and mature aligns the recognition task with the actual decisions smallholder farmers must take, especially identifying the harvest-ready state. Because phase transitions are continuous and visually subtle, an on-demand automated system gives a more consistent assessment than periodic manual inspection. The dominance of sub-hectare landholdings also requires the monitoring solution to be low in hardware and operating cost, and free from specialized expertise.

\section{Proximal Imaging}
\label{sec:proximal_imaging}

Crop sensing technologies are grouped by the proximity of the sensor to the canopy. Remote sensing platforms such as satellite constellations and Unmanned Aerial Vehicles (UAVs) observe fields from tens of meters to several hundred kilometers with spatial resolutions from ten meters down to centimeters, but at the cost of platform mobilization expense and dependence on atmospheric conditions. Proximal sensing places sensors within centimeters to a few meters of the plants and provides real-time, ground-level measurements without those dependencies \cite{mezera2021proximalremote}. A direct comparison in winter wheat across four growing seasons showed that proximal on-the-go sensors remain operable under cloud cover that prevents reliable satellite acquisition, a recurring constraint during critical phenological transitions \cite{mezera2021proximalremote}.

For rice, proximal imaging with consumer-grade Red-Green-Blue (RGB) cameras is a low-cost alternative to multispectral remote sensing platforms that remain underutilized in smallholder contexts. A study covering more than 4{,}800 harvesting plots across six countries showed that vertically downward RGB images of the rice canopy at 0.8 to 0.9 meter, resolved over one square meter consistent with FAO harvesting standards, support accurate yield estimation without labor-intensive crop cuts \cite{tanaka2023rgbrice}. The same study framed proximal RGB imaging as a practical pathway for rice monitoring in the global South, where smallholder land structure makes large-area remote sensing economically unattractive.

This study follows the same proximal RGB principle. A single consumer-grade RGB camera is mounted statically above the rice canopy for close-range, top-down, on-demand capture, rather than continuous video or oblique handheld photography. The static top-down arrangement constrains the geometric variability of the field of view, and the on-demand trigger makes bandwidth, storage, and inference cost scale with farmer interaction rather than sensor duty cycle. The setup trades spatial coverage of UAV or satellite sensing for plot-level temporal flexibility and operational simplicity.

Proximal RGB imaging in the field is not without challenges. Cameras must operate under uncontrolled illumination, and ground-based rice yield studies report that imaging time of day measurably affects color balance and contrast \cite{tanaka2023rgbrice}. Motion blur from mounting vibration, wind, or platform displacement degrades downstream vision tasks, with controlled experiments showing segmentation accuracy on crop and weed images dropping from a mean intersection-over-union of about 0.80 on sharp images to about 0.63 on motion-blurred ones \cite{yun2023motiondeblurring}. The camera may also capture frames containing soil, equipment, water, or other non-rice content that are uninformative for phenology classification.

A closed-world classifier would assign a phenology label to images that contain no rice canopy or that are otherwise unusable for phenology interpretation, so the inference pipeline carries an open-world rejection step at its entry. A Vision-Language Model gatekeeper enforces this rejection before the supervised classifier receives the frame.

\begin{figure}
    \centering
    \includegraphics[width=0.55\textwidth]{img/Overhead canopy capture.png}
    \caption{Overhead RGB Canopy Capture}
    \label{fig:proximal-imaging}
\end{figure}

\section{ESP32-CAM}

The ESP32-CAM is a compact, low-cost embedded vision module built around the Espressif ESP32 system-on-chip and a camera daughterboard. The ESP32 integrates a dual-core 32-bit microcontroller, several hundred kilobytes of on-chip Static Random Access Memory (SRAM), and integrated 2.4 GHz Wi-Fi and Bluetooth connectivity, which makes it a popular choice for Internet of Things deployments that combine wireless data transport with on-device control \cite{hercog2023esp32iot}. In the AI-Thinker ESP32-CAM variant used in this study, the system-on-chip is paired with 4 MB of external Pseudo Static Random Access Memory (PSRAM) to buffer image frames, 4 MB of Serial Peripheral Interface (SPI) flash for firmware, and a microSD card slot for optional local logging. The full imaging node fits within roughly 27 by 40 millimeters at a unit cost well below ten United States dollars, which is the order of magnitude required for a monitoring solution that smallholder farmers can replicate.

The optical front end is the OmniVision OV3660, a three-megapixel Complementary Metal Oxide Semiconductor (CMOS) sensor that supports capture up to 2048 by 1536 pixels with on-sensor Joint Photographic Experts Group (JPEG) compression. Compared with the OV2640 commonly fitted to default AI-Thinker boards, the OV3660 offers higher native spatial resolution, which is useful for resolving rice canopy texture and panicle morphology at close range. The sensor is driven by the ESP32 over a parallel data interface and an Inter-Integrated Circuit (I2C) control bus, and JPEG-encoded frames are buffered in PSRAM before transmission, so no companion processor is required on the acquisition path.

The ESP32 platform has been adopted as the imaging and sensing endpoint of several plant monitoring systems. Elhattab et al. proposed a remote plant growth monitoring system in which an ESP32-CAM transmits images over Wi-Fi to a mobile application, with the server side handling visualization and growth tracking \cite{elhattab2023esp32cam}. Correa-Quiroz et al. reported a greenhouse system in which an ESP32 microcontroller aggregates climate, soil moisture, and irrigation telemetry and streams them to a cloud dashboard \cite{correaquiroz2025esp32greenhouse}. The same offloading pattern recurs outside agriculture in beehive and fermentation monitoring nodes that forward local measurements to an Internet of Things cloud platform such as ThingSpeak \cite{hercog2023esp32iot}. Across these deployments the ESP32 acts as a capture-and-upload endpoint rather than an inference platform, with analytical workload deferred to a more capable backend. The few hundred kilobytes of SRAM and few megabytes of PSRAM cannot accommodate the parameter footprint of mainstream Convolutional Neural Networks (CNNs) at full precision, and even quantized models leave little headroom for the runtime, the input pipeline, and the JPEG encoder.

In this study, the ESP32-CAM is configured strictly as a periodic image acquisition node. On each wake from a deep-sleep cycle, the device captures a single JPEG frame at full resolution, attaches a timestamp and device identifier, and uploads the frame to cloud storage over Wi-Fi before returning to sleep. The device performs no preprocessing beyond the sensor's built-in pipeline and does not run any classification model locally. This division keeps the embedded firmware small and predictable, concentrates the energy budget on short Wi-Fi bursts rather than continuous inference, and allows classifier updates to be rolled out by changing the cloud service alone.

The hardware choice introduces constraints. The OV3660 lacks the radiometric calibration interfaces of dedicated agricultural cameras, and its automatic exposure and white balance operate on a single integrated RGB Bayer sensor, so any colorimetric analysis downstream is implicitly conditioned on the image signal processor. The static top-down mounting reduces but does not eliminate pose variability, and the on-demand capture mode reduces motion blur compared with continuous video, although ambient illumination changes during the day remain uncontrolled. These constraints are accepted as part of a trade-off that prioritizes hardware affordability over laboratory-grade radiometric fidelity.

\begin{figure}
    \centering
    \includegraphics[width=0.5\textwidth]{img/espcam.png}
    \caption{ESP32-CAM}
    \label{fig:esp32cam}
\end{figure}

\section{Image Features}

In addition to raw pixel intensities, this study uses compact image features derived from each RGB frame as inputs to the classical machine learning classifiers and to the handcrafted branch of the hybrid model. Two families are used in parallel: color vegetation indices from the per-pixel red, green, and blue channels, and Gray Level Co-occurrence Matrix (GLCM) texture descriptors from a grayscale projection of the frame. Color vegetation indices capture chromatic information sensitive to chlorophyll content and canopy greenness, while GLCM descriptors capture the spatial arrangement of intensities that reflects canopy structure. Recent work on cotton yield estimation showed that fusing texture and color features from an RGB sensor outperforms either family alone, raising the coefficient of determination of the yield model to 0.91 on validation data \cite{ma2022cottonRGB}. A similar pattern has been reported for rice canopy estimation from Unmanned Aerial Vehicle (UAV) imagery, where texture features compensate for spectral information at early stages dominated by soil and water background noise and at later stages affected by spectral saturation as the canopy closes \cite{liu2023glcmrice}.

\subsection{Color Vegetation Indices}

Color vegetation indices are analytical combinations of the red, green, and blue channels that emphasize green vegetation against soil, residue, or shadow. They serve as low-cost substitutes for multispectral indices such as the Normalized Difference Vegetation Index (NDVI) when only a consumer-grade RGB camera is available, since they can be computed without a dedicated near-infrared band. Recent RGB-based studies report strong correlations between selected color indices and crop biophysical traits across cereal, cotton, and cover crop systems \cite{ma2022cottonRGB}.

This study uses eight color vegetation indices widely reported in RGB-based precision agriculture: the Excess Green (ExG), Excess Red (ExR), Excess Green minus Excess Red (ExGR), Visible Atmospherically Resistant Index (VARI), Green Leaf Index (GLI), Normalized Green Red Difference Index (NGRDI), Modified Green Red Vegetation Index (MGRVI), and Red Green Blue Vegetation Index (RGBVI) \cite{ma2022cottonRGB}. Let $R$, $G$, and $B$ denote the raw red, green, and blue intensities of a pixel. The eight indices are defined as:

\begin{figure}
    \centering
    \includegraphics[width=0.85\textwidth]{img/Color Vegetation Indices.png}
    \caption{Color Vegetation Indices}
    \label{fig:vi-indices}
\end{figure}

The eight indices fall into two groups. ExG, ExR, and ExGR are difference-style indices that compare a single emphasized channel against the others. ExGR is reported as an effective separator of vegetation from soil and senescent residue because the excess red term suppresses green-red material such as stems and dry tissue that ExG alone tends to misclassify as foliage. NGRDI and MGRVI are ratio-style indices that normalize a green-red contrast against the green-plus-red sum, while GLI normalizes a green-versus-non-green contrast against the full red-green-blue sum, and RGBVI compares the squared green channel against the red-blue product. VARI substitutes blue in the denominator with the opposite sign $(G-R)/(G+R-B)$, which is the atmospheric-resistance correction that gives it robustness against scene-wide haze and brightness changes \cite{ma2022cottonRGB}. Across the rice phenological cycle, the relative proportions of green foliage, panicle, and senescent tissue evolve systematically, so both groups contribute discriminative information for distinguishing the vegetative, generative, and mature phases.

For each RGB frame, the eight indices are computed pixel-wise, and ten descriptive statistics are extracted from each per-index map to form a fixed-length feature vector summarizing the chromatic state of the canopy.

\subsection{GLCM Texture Features}

Color vegetation indices summarize chromatic information at the pixel level but discard the spatial arrangement of intensities. Texture descriptors recover this information by analyzing the statistical relationship between neighboring pixel intensities. The Gray Level Co-occurrence Matrix (GLCM) is a second-order statistical representation that counts the joint occurrence of pixel intensity pairs at a given spatial offset, parameterized by a displacement distance and an angular direction. From a normalized GLCM, scalar descriptors can be derived that capture distinct aspects of local intensity variation.

This study uses four GLCM descriptors widely adopted in GLCM-based crop classification from low-altitude remote sensing imagery: Contrast, Correlation, Energy (also called Angular Second Moment), and Homogeneity \cite{iqbal2021glcmcrop}. Contrast responds to sharp intensity transitions and highlights edge density, while Homogeneity responds inversely and highlights uniform regions. Energy quantifies the textural orderliness of the co-occurrence distribution and increases when the canopy texture is regular. Correlation measures how predictable a neighboring pixel intensity is given the reference pixel and responds to directional structure in the texture. Each descriptor is computed at two pixel offsets of one and three pixels and four angular directions of $0$, $\pi/4$, $\pi/2$, and $3\pi/4$ radians, yielding eight cells per descriptor. The mean and standard deviation across these eight cells are then taken as two summary statistics for that descriptor, so the four descriptors together produce eight texture features that capture both the average textural state and its directional and scale anisotropy.

For the classical machine learning experiments, the eight texture features are concatenated with the ten per-index statistics from each color vegetation index, giving a handcrafted feature vector of eighteen dimensions per vegetation index. The same vector is reused as the handcrafted branch of the hybrid model. Combining texture features with chromatic information is consistent with recent findings that texture compensates for spectral information at early growth stages affected by background interference and at late stages affected by spectral saturation \cite{liu2023glcmrice,ma2022cottonRGB}.

\section{Machine Learning}

The classical machine learning track of this study applies supervised classifiers to the eighteen-dimensional handcrafted feature vector formed by color vegetation index statistics and Gray Level Co-occurrence Matrix descriptors. The motivation is operational: classical classifiers have small parameter counts and low inference cost, which suits a cloud backend that charges per request, and their feature attributions are easier to communicate to non-technical stakeholders. A recent review of machine learning in agriculture confirms that Random Forest, Support Vector Machine, K-Nearest Neighbors, Logistic Regression, and Gradient Boosting remain widely used for crop monitoring, with each method offering a different inductive bias \cite{benos2021ml}. For paddy in particular, Random Forest has been reported to outperform Classification and Regression Trees (CART) and Support Vector Machine on growth phase classification with imbalanced data \cite{suryono2022paddyrf}, the same regime present in the three-class vegetative, generative, and mature partition used here. Five classifiers are evaluated, each covering a distinct inductive family.

\subsection{Random Forest}

Random Forest (RF) is an ensemble of decision trees trained on bootstrap samples, with each split chosen from a random subset of features, and predictions aggregated by majority vote. The construction reduces the variance of a single decision tree, captures non-linear feature interactions, and tolerates heterogeneous feature scales and moderately many irrelevant features. These properties suit the mixed vegetation index statistics and texture descriptors used here, and Random Forest has been reported as a strong classifier for paddy growth phase classification under class imbalance \cite{suryono2022paddyrf}.

\subsection{Support Vector Machine}

Support Vector Machine with a Radial Basis Function kernel (SVM-RBF) constructs a maximum-margin separator in a feature space defined implicitly by a Gaussian kernel, which allows non-linear class boundaries without explicit polynomial expansion. SVM-RBF fits moderate-dimensional tabular vectors well when the sample count is comparable to or larger than the feature dimension, and it is a common baseline in supervised crop monitoring studies \cite{benos2021ml,suryono2022paddyrf}.

\subsection{K-Nearest Neighbors}

K-Nearest Neighbors (KNN) assigns a query point the majority class of its $k$ closest training points under a chosen distance metric, with no parametric training stage. KNN provides a non-parametric, instance-based baseline that is sensitive to the local geometry of the feature distribution and indicates whether the handcrafted feature space admits a simple nearest-neighbor structure aligned with the phenological classes \cite{benos2021ml}.

\subsection{Logistic Regression}

Logistic Regression (LR) fits a linear decision boundary in the input feature space and produces class probabilities through the softmax extension for multiclass problems. As a linear baseline, it indicates whether the handcrafted feature vector is linearly separable across the three phases, which is informative about the geometric difficulty of the task \cite{benos2021ml}.

\subsection{Gradient Boosting}

Gradient Boosting (GB) builds an additive ensemble of shallow decision trees by sequentially fitting each new learner to the residuals of the current ensemble. The stage-wise optimization concentrates capacity on samples that are still misclassified, which is effective on heterogeneous tabular features such as the combined vegetation index and texture descriptors used here \cite{benos2021ml}.

The five classifiers are trained and evaluated on the handcrafted feature vector derived from each of the eight color vegetation indices, producing forty combinations of vegetation index and classifier. All combinations share the same cross-validation protocol so that the contribution of feature choice and the contribution of classifier choice can be isolated.

\section{Deep Learning}

The deep learning track of this study replaces the handcrafted feature pipeline with Convolutional Neural Networks (CNNs) that learn a hierarchy of spatial filters directly from raw pixels. Unlike the vegetation index and Gray Level Co-occurrence Matrix descriptors used in the classical track, CNN filters are trained jointly with the classifier head, so chromatic and structural information are extracted by the same network and optimized for the phenology objective. The cost is a larger parameter count and a stronger dependence on training data volume, which is addressed by transfer learning from networks pretrained on the ImageNet classification dataset. The pretrained backbones provide general-purpose low and mid-level visual features that can be reused for plant imagery after lightweight adaptation, which is the established practice across plant image recognition on small in-field datasets, including rice growth stage classification \cite{pacal2024dlplantdisease,qin2023ricephenology}.

CNN-based recognition of rice growth stages has been demonstrated at both the canopy and the panicle level. Qin et al. trained a ResNet-50 backbone on a six-stage rice phenology dataset of RGB canopy imagery captured across the full life cycle and reported around 87 percent recognition accuracy across the six stages \cite{qin2023ricephenology}. Tan et al. introduced the RiceRes2Net architecture for in-field panicle detection and growth stage recognition and showed that a residual CNN backbone tuned on rice imagery captures fine-grained morphological cues such as panicle emergence and grain filling that are hard to encode by hand \cite{tan2023riceres2net}. These studies establish CNN backbones as a competitive choice for in-field rice phenology recognition and motivate the three backbones evaluated here, selected to span the range from accuracy-oriented designs down to ultra-compact mobile-oriented designs at progressively smaller parameter budgets.

All three backbones are fine-tuned on the curated rice dataset using a two-phase transfer learning protocol. The classification head is first trained with the backbone frozen so that the new softmax weights adapt to the rice feature distribution without disturbing the pretrained filters. The top layers of the backbone are then unfrozen and trained jointly with the head at a lower learning rate to refine high-level features for the target task while preserving the generic low-level filters learned on ImageNet. The same protocol is reused for the visual branch of the hybrid model so that the deep learning and hybrid tracks share their feature extractor under identical training conditions.

\subsection{EfficientNetB0}

EfficientNetB0 is the baseline of the EfficientNet family of CNNs designed through compound scaling, a principled scheme that jointly scales network depth, width, and input resolution under a fixed computational budget \cite{tanle2019efficientnet}. The backbone is built from mobile inverted bottleneck convolutional blocks with squeeze-and-excitation modules and produces a 1280-dimensional embedding from a global average pooling head before the final classification layer \cite{tanle2019efficientnet}. The architecture reaches high ImageNet accuracy with a parameter count of about 5.3 million, which is attractive when classification accuracy is prioritized but the per-request inference cost of a cloud backend still has to remain modest.

In this study, EfficientNetB0 is the primary deep learning backbone for the rice phenology classifier. The ImageNet pretrained weights are reused, the original classification head is replaced by a three-class softmax layer matching the vegetative, generative, and mature labels, and the backbone is fine-tuned on the curated rice dataset. The same backbone is also the visual branch of the hybrid model, so its 1280-dimensional embedding is shared between the deep learning track and the hybrid track and the two pipelines remain directly comparable.

\begin{figure}
    \centering
    \includegraphics[width=0.85\textwidth]{img/EfficientNet-Architecture-diagram.jpg}
    \caption{EfficientNet Architecture}
    \label{fig:effnet-architecture}
\end{figure}

\subsection{MobileNetV2}

MobileNetV2 is a compact CNN backbone originally designed for mobile and embedded deployment, although it is used here as a server-side classifier \cite{sandler2018mobilenetv2}. The architecture is built from inverted residual blocks with linear bottlenecks, in which channel-wise depthwise separable convolutions replace dense spatial convolutions and residual shortcuts connect the low-dimensional bottleneck representations rather than the expanded feature maps \cite{sandler2018mobilenetv2}. The combined effect is a substantial reduction in parameter count, on the order of 3.5 million for the standard configuration, while preserving competitive ImageNet accuracy.

MobileNetV2 enters this study as a lightweight comparison backbone for EfficientNetB0. The same transfer learning protocol is applied, with ImageNet pretrained weights loaded, the classification head replaced by a three-class softmax, and the network fine-tuned on the rice dataset. Its smaller parameter footprint indicates whether the additional capacity of EfficientNetB0 is justified for the three-phase rice phenology task and how much accuracy can be retained when the backbone is sized down toward mobile deployment.

\subsection{MobileNetV3-Small}

MobileNetV3-Small is the most compact configuration in the MobileNetV3 family, designed for execution on devices with very limited compute and memory \cite{howard2019mobilenetv3}. The architecture is the product of Neural Architecture Search combined with the NetAdapt algorithm, and it introduces two changes that distinguish it from the previous MobileNet generation: the h-swish activation, a hardware-friendly approximation of the swish function, and the integration of squeeze-and-excitation modules in selected blocks \cite{howard2019mobilenetv3}. The Small variant carries about 2.5 million parameters and targets latency budgets at which even MobileNetV2 is too costly.

In this study, MobileNetV3-Small is the lightest backbone evaluated and serves as the lower bound on backbone capacity for rice phenology classification. MobileNetV3 has been adopted as a lightweight backbone in related rice agricultural tasks, including pest and disease detection in field imagery, where its low parameter count supports deployment scenarios that EfficientNet-scale models cannot reach \cite{jia2023ricemobilenetv3}. The same two-phase transfer learning protocol used for the other backbones is applied, with ImageNet pretrained weights loaded and the three-class softmax head replacing the original classifier. Inference still runs in the cloud backend, but the small parameter footprint of MobileNetV3-Small benchmarks the lower end of the accuracy versus capacity trade-off relevant for future edge deployment at the sensor node.

\section{Hybrid Model}

The hybrid track combines the learned visual features of a Convolutional Neural Network (CNN) with the handcrafted vegetation index and texture features used in the classical track within a single classifier. The motivation is complementarity: the CNN backbone captures spatial structure and high-level abstraction that pixel-level color statistics and Gray Level Co-occurrence Matrix (GLCM) descriptors cannot represent, while the handcrafted features inject explicit chromatic and textural priors that are aligned with rice canopy biology and were already shown to be informative in the classical track. Recent crop-phenology work has validated that late fusion of CNN-extracted features with handcrafted descriptors improves fine-grained phenological recognition over either family alone, reaching accuracies above 0.92 on crop-specific phenology benchmarks \cite{yang2024phenologynet}.

The architecture used in this study is a late-fusion two-branch network. The visual branch is an EfficientNetB0 backbone with ImageNet pretrained weights and the original classification head removed; its output is reduced by global average pooling to a 1280-dimensional embedding that summarizes the canopy at a high semantic level. The handcrafted branch receives the 18-dimensional feature vector formed by 10 statistics of a single vegetation index map and 8 Gray Level Co-occurrence Matrix texture features computed from the grayscale projection of the same frame, and applies batch normalization to bring the two branches onto comparable scales. The two embeddings are concatenated and passed through a small classification head of two fully connected layers with 128 and 64 ReLU units, separated by a dropout layer with rate 0.3, before a final three-class softmax produces the vegetative, generative, or mature prediction.

Training follows the same two-phase transfer learning protocol used for the standalone CNN track. The classification head and the batch normalization on the handcrafted branch are first trained for ten epochs with the backbone frozen at a learning rate of $10^{-3}$, so that the new head adapts to the joint representation without disturbing the pretrained filters. The top thirty layers of the backbone are then unfrozen and the full network is fine-tuned for up to twenty-five epochs at a reduced learning rate of $10^{-5}$, with early stopping on validation Macro F1 at a patience of five epochs. Class weights are computed inversely to class frequency in the training partition so that the underrepresented mature class receives sufficient gradient signal during both phases.

Only three vegetation indices are evaluated in the hybrid track: ExGR, VARI, and NGRDI, selected as the three best-performing indices identified in the classical machine learning track. This restriction holds the classifier capacity constant and isolates the effect of the chromatic prior on hybrid performance, so that the contribution of vegetation index choice can be compared against the same choice made within a non-deep classifier. The result is three hybrid models per fold that share architecture, training schedule, and image branch, and that differ only in the 18-dimensional handcrafted vector concatenated with the CNN embedding.

\section{Vision-Language Model}

A Vision-Language Model (VLM) is a multimodal large language model that accepts both visual and textual inputs and produces textual outputs grounded in the image content. Compared with a dedicated image classifier that learns a fixed mapping from pixels to a closed label set, a VLM is conditioned at inference time by a natural language instruction, so the label set, the output format, and the decision criteria can be redefined per request without retraining. A recent survey of vision-language models for vision tasks organizes the family along three axes, pretraining objective, transfer learning paradigm, and prompt design, and documents that VLMs trained on large image-text corpora transfer to novel visual recognition tasks under zero-shot prompting without any task-specific fine-tuning \cite{zhang2024vlmsurvey}.

Recent agricultural studies have begun to evaluate whether VLMs can directly substitute supervised classifiers on field imagery. A 2026 benchmark of six commercial multimodal models, including ChatGPT-4o and Gemini Flash 2.5, evaluated zero-shot weed detection in Unmanned Aerial Vehicle (UAV) imagery and reported that the choice of prompt and the inclusion of counterfactual probing materially shape both classification accuracy and the interpretability of the model's reasoning \cite{frontiers2026weedvlm}. The same study positioned commercial VLMs as low-annotation tools that lift the requirement for large labeled training sets, while emphasizing that prompt construction, not raw model capacity alone, governs reliable field-grade performance \cite{frontiers2026weedvlm}.

A closed-set classifier assumes that every input belongs to one of the training classes, an assumption that breaks down once a field-deployed camera can capture frames containing soil, equipment, water, motion blur, or any non-rice content. Open-set recognition formalizes this gap by treating unknown classes as a separate decision outcome and reducing the open space risk incurred by force-fitting them into known categories, and deep learning extensions of the same idea replace the standard softmax layer with a calibrated decision rule that explicitly models unknowns and rejects them at inference time \cite{geng2021openset}. Without an open-set safeguard, a phenology classifier would assign a vegetative, generative, or mature label to any frame regardless of its rice content, which is unsafe for an on-demand farmer-facing system.

In this study, the VLM is used in two related roles. As an open-world gatekeeper, the VLM receives each uploaded frame before the closed-set phenology classifier and is prompted to decide whether the frame contains a rice canopy at adequate quality, returning a binary in-distribution or out-of-distribution outcome. As a candidate phenology classifier in the comparative track, the same VLM is also prompted to return one of the three phenological labels directly from the image, without any supervised fine-tuning on the curated rice dataset. Both roles share the same model and provider account, which keeps deployment surface and billing footprint identical between gatekeeping and classification.

Four proprietary VLMs from two providers, Google and OpenAI, are evaluated in the comparative track. Open-source VLMs are excluded from the present scope because their self-hosted deployment cost and engineering effort do not match the cloud-only operating model targeted here.

\subsection{Gemini 2.5}

The Gemini family of multimodal models, developed by Google, is a frontier vision-language model line designed for joint understanding of text, images, audio, and video, and recent surveys of vision-language models catalogue it alongside the GPT-4 family as a reference point in the current generation of large multimodal systems \cite{zhang2024vlmsurvey}. Two endpoints from the 2.5 generation are evaluated for a price-sensitive cloud deployment: Gemini 2.5 Flash, the mid-tier balance of latency and accuracy, and Gemini 2.5 Flash Lite, the most compact endpoint optimized for high request volume at the lowest per-call cost in the 2.5 line.

The two Gemini 2.5 endpoints are queried through the official multimodal Application Programming Interface (API) with identical prompt content. Holding prompt and protocol constant across the two endpoints isolates the effect of capacity tier within a single provider on zero-shot rice phenology recognition.

\subsection{GPT-4o}

GPT-4o is the multimodal extension of the GPT-4 family released by OpenAI, with native support for image inputs and instruction following at the prompt level, and it is catalogued among the frontier proprietary vision-language models in recent surveys of the field \cite{zhang2024vlmsurvey}. Two endpoints are evaluated: GPT-4o, the full-capacity model targeted at general multimodal reasoning, and GPT-4o mini, a smaller variant designed for high-throughput deployment at reduced latency and lower per-call cost. The two endpoints span the same accuracy versus cost axis as the Flash and Flash Lite pair on the Gemini side, which allows cross-provider comparison at comparable price tiers rather than only at the top capacity tier.

Both endpoints are queried with the prompt template and image preprocessing pipeline used for the Gemini endpoints. The resulting four-model comparison isolates the effect of provider choice and capacity tier on zero-shot rice phenology recognition under the same dataset and evaluation protocol applied to the classical, deep, and hybrid tracks.

\section{Evaluation Metrics}

Every pipeline track in this study, classical, deep, hybrid, and vision-language, must be evaluated under the same protocol so that performance differences reflect modeling choices rather than evaluation variability. The metrics adopted here are the standard supervised classification family derived from the confusion matrix, extended to the multi-class setting through macro averaging, and accompanied by a resampling-based protocol for both training-time selection and post hoc uncertainty quantification. The same suite is widely adopted in agricultural machine learning studies, where class imbalance, small sample counts, and the operational cost of misclassification motivate reporting more than a single scalar accuracy \cite{benos2021ml}.

\subsection{Confusion Matrix}

The confusion matrix is a square table whose rows index the ground-truth class and whose columns index the predicted class, with each cell counting the number of samples mapped from row class to column class. For a $K$-class problem the matrix has $K \times K$ cells and exposes the full pattern of correct predictions on the diagonal and class-pair-specific errors off the diagonal. In the three-class rice phenology setting used here, with classes vegetative, generative, and mature, the confusion matrix isolates errors that matter operationally, such as predicting mature when the canopy is still generative, from errors that matter less, such as confusion between phases that share visual cues.

\subsection{Accuracy}

Accuracy is the fraction of samples whose predicted class equals the ground-truth class, computed as the sum of the diagonal of the confusion matrix divided by the total sample count. Accuracy summarizes overall correctness in a single scalar but is dominated by the majority class when class frequencies are uneven, which means it can understate the cost of failing on a minority class. The mature class in this study is operationally the most critical class but is not necessarily the largest in the training partition, so accuracy is reported but is not used as the primary selection criterion.

\subsection{Precision}

Precision for a class is the fraction of samples predicted as that class that are actually members of it, defined as
\begin{equation}
    \mathrm{Precision} = \frac{\mathrm{TP}}{\mathrm{TP} + \mathrm{FP}}
    \label{eq:precision}
\end{equation}
where $\mathrm{TP}$ and $\mathrm{FP}$ are the per-class true positive and false positive counts read from the confusion matrix. Precision penalizes false alarms and is directly relevant to the cost of a system telling a smallholder farmer that the crop is mature when it is not, since a premature harvest decision based on a false positive degrades both yield and grain quality. Per-class precision is computed for each of the three rice phenology classes.

\subsection{Recall}

Recall for a class is the fraction of actual members of that class that the classifier successfully predicts as such, defined as
\begin{equation}
    \mathrm{Recall} = \frac{\mathrm{TP}}{\mathrm{TP} + \mathrm{FN}}
    \label{eq:recall}
\end{equation}
where $\mathrm{FN}$ is the per-class false negative count. Recall penalizes missed detections and is directly relevant to the cost of failing to flag a mature canopy that the farmer should harvest, since a missed positive at the mature phase delays harvest and increases shattering losses. Recall on the mature class is therefore tracked as the IoT-critical secondary metric of this study, alongside per-class recall for vegetative and generative.

\subsection{F1-Score}

The F1-score for a class is the harmonic mean of precision and recall for that class, defined as
\begin{equation}
    \mathrm{F1} = \frac{2 \cdot \mathrm{Precision} \cdot \mathrm{Recall}}{\mathrm{Precision} + \mathrm{Recall}}
    \label{eq:f1}
\end{equation}
The harmonic mean penalizes large imbalance between precision and recall, so a class with high recall but very low precision produces a low F1-score, and conversely. For multi-class evaluation, the per-class F1 scores are aggregated by macro averaging, defined as the unweighted arithmetic mean across the $K$ class-specific F1 scores. Macro averaging treats each class as equally important regardless of its size, which is appropriate when minority classes carry equal or higher operational weight than majority ones, and the statistical inference for the macro-averaged F1 score has been formalized through a recent confidence interval derivation that explicitly handles the multi-class extension \cite{takahashi2022f1ci}. The Macro F1 score is the primary model-selection metric of this study.

\subsection{Cross-Validation}

Cross-validation is a resampling protocol that uses the labeled data set itself to estimate generalization performance, by partitioning the data into $k$ disjoint folds, training on $k-1$ folds, and evaluating on the held-out fold, then rotating until each fold has served once as the held-out partition. Stratified $k$-fold cross-validation enforces that each fold preserves the class-frequency distribution of the original data, which mitigates evaluation variance under class imbalance and avoids degenerate folds that contain no samples of a minority class \cite{szeghalmy2023stratifiedcv}. This study uses 5-fold stratified cross-validation on the 1036-image training partition for every candidate model, with the per-fold Macro F1 score averaged to produce the cross-validated estimate used for ranking across the classical, deep, hybrid, and vision-language tracks.

\subsection{Bootstrap Confidence Interval}

Single-point estimates of classification metrics do not communicate their sampling uncertainty, which is essential when small differences in cross-validated Macro F1 must be compared across many candidate models. The bootstrap procedure resamples the test set with replacement to construct an empirical sampling distribution of the metric of interest, from which a percentile-based confidence interval can be read directly. Analytical confidence intervals for macro-averaged and micro-averaged F1 scores have been derived under large-sample asymptotics using a multivariate central limit theorem \cite{takahashi2022f1ci}, while the non-parametric bootstrap is the practical alternative when the test sample is small and a single resampling procedure must apply uniformly to several metrics. The bootstrap protocol is applied here to the held-out test set of 104 images to attach a 95\% confidence interval to every reported Macro F1, Mature Recall, and per-class score, so that model rankings can be interpreted with their uncertainty bands rather than treated as exact differences.

\section{Cloud IoT Backend}

The Internet of Things (IoT) places sensing devices, network connectivity, and cloud-side processing into one system. A comprehensive review of IoT in smart farming identifies cloud computing as the dominant analytical and storage layer because field-side compute budgets are typically too constrained for inference and full-cycle retention \cite{boursianis2022iot}. A canonical cloud-IoT pipeline for smart agriculture follows a three-layer pattern of perception, transport, and cloud, with the cloud center performing deep learning while perception devices remain low-cost and stateless \cite{yu2025cloudedge}. The same line of work quantifies the latency and bandwidth trade-off between local and cloud processing for image classification in smart agriculture and shows that cloud execution is favored when inference workloads are sparse and model updates are frequent \cite{markovic2024cloudfog}.

The backend used in this study realizes the three-layer pattern with cloud-only inference. ESP32-CAM nodes produce JPEG frames on a periodic deep-sleep cycle and upload them over Wi-Fi to a Google Cloud Storage (GCS) bucket. A Flask application on Google Cloud Run loads each frame, runs the selected classifier from the comparative pool, and emits the predicted phenology class with its probability vector. Operational state, including the latest predicted class, confidence, model identifier, and capture timestamp, is written to a Firebase Realtime Database (RTDB) that the dashboard subscribes to. The dashboard is a Progressive Web Application (PWA) hosted on Netlify, accessible from any smartphone.

Two consequences follow. The perception layer is stateless, so a backend update reaches every deployed unit without firmware re-flashing, and the inference compute is billed per request through the serverless Cloud Run runtime, which removes idle-time cost \cite{yu2025cloudedge}. The accepted costs are a network round trip plus a cold-start overhead of several hundred milliseconds on a scaled-to-zero instance, and a hard dependence on Wi-Fi at the perception node. Both fit the same affordability-over-fidelity trade-off that motivates the ESP32-CAM hardware choice.

\section{Regulation and Data Governance}

A cloud-based IoT system for smallholder agriculture handles two governance regimes at once. Canopy images and operational telemetry from the ESP32-CAM are non-personal in narrow terms but carry implicit information about a specific farmer's plot, while farmer-side identifiers and access credentials are personal data in the strict legal sense. Recent governance research catalogues smallholder-side privacy risks including farmer identification, reputational harm, and limited downstream control once data has been uploaded \cite{kaur2022farmerprivacy}. A systematic review of agricultural data ownership further distinguishes legal, voluntary, and technical enforcement as parallel pathways and notes that much non-personal agricultural data falls outside personal data protection legislation, which leaves the ownership question to contractual and technical mechanisms \cite{berisha2026dataownership}.

International IoT standardization sets the architectural baseline. ITU-T Recommendation Y.2060 lists interconnectivity, autonomic data exchange, and embedded security as defining IoT characteristics and treats application-layer integrity and authentication as design obligations \cite{itut_y2060}. ISO/IEC 30141:2024 decomposes a conformant IoT system into perception, network, application, and security layers and requires that data collection, processing, transmission, and protection be addressed at every layer rather than aggregated into a single control point \cite{iso30141_2024}.

Indonesia's electronic information framework attaches binding legal weight to the same layers. Law Number 19 of 2016 amending Law Number 11 of 2008 classifies any digitally encoded information, including the JPEG frames and prediction outputs handled by the backend, as electronic information whose integrity and authenticity must be protected against unauthorized access \cite{uuite2016}. Government Regulation Number 71 of 2019 imposes corresponding obligations on the electronic system operator role occupied by the cloud backend, including security controls, reliability, and the suitability of hardware and software \cite{pp71_2019}.

Personal data is governed by a more recent regime. Law Number 27 of 2022 on Personal Data Protection establishes the first comprehensive personal data protection law in Indonesia and defines personal data as information about an identified or identifiable natural person, including identifiers combinable with other information to single out an individual \cite{uupdp2022}. For a smallholder-facing platform, the relevant attack surface is the linkage between farmer identifiers in the Realtime Database and image records in cloud storage, since the joint dataset can re-identify a plot owner even when the canopy images themselves contain no personal features.

The translation of these regulatory positions into concrete data governance controls, including identifier separation, storage residency, retention lifecycle, and audit logging, is detailed in Chapter~3.
